\documentclass[11pt,a4paper]{article}
\usepackage[T1]{fontenc}
\usepackage{lmodern,amsmath,amssymb,amsthm,booktabs,microtype,hyperref}
\usepackage[margin=27mm]{geometry}
\setlength{\emergencystretch}{3em}
\newcommand{\Q}{\mathbb Q}\newcommand{\Z}{\mathbb Z}
\newcommand{\Tr}{\operatorname{Tr}}\newcommand{\Nm}{\operatorname{N}}
\newtheorem{lemma}{Lemma}
\newtheorem{proposition}[lemma]{Proposition}
\newtheorem{theorem}[lemma]{Theorem}
\title{Prepared native products with no adaptive square witness}
\author{}\date{17 November 2026}
\begin{document}\maketitle
\begin{abstract}
Two admissible adaptive square classes fail to have a rational
witness at some fully prepared native detector inputs, even when
their added scalar is unrestricted. The first class is the previously
retained quartic in $q=2+wc^2$. The second is $1+w^2V^4c^4$.
For each, an exact positive native product over $\Q(\sqrt2)$ gives
a pole input with square reduced denominator and all intrinsic-pole
preparation conditions, but a local obstruction at $100003$ excludes
every scalar value. At any fixed finite set of primes away from
the obstruction prime, arbitrary prescribed finite-place and real
closeness to the reference product can also be imposed. The second
square equation is identified with a quartic covering of a prescribed
congruent-number curve. These are abstract native products, not
verified actual elliptic source tuples; the results do not disprove
the twenty-one- or nine-variable realization criteria.
\end{abstract}

\section{Two scalar criteria and their scope}

Use $K=\Q(\sqrt2)$ and write
\[
 w=\frac{u}{(1+u)^2},\qquad
 W=\prod_i(1+wh_i^2)\in K^\times,\qquad
 V=\Tr(W^{-1}),\qquad b=\frac{4w}{V^2},
\]
with $h_i\in K$ and all displayed divisions defined. Set
\[
                         C=V^2+w(2-V).
\]
The trace criterion of \cite{Trace} says that rational squareness
of $C$ excludes the negative odd poles used by the source, when
every odd valuation of $w$ is positive. Its adaptive class is
\begin{equation}\label{eq:first}
 q=2+wc^2,\qquad
 \Delta(q,V)=(q^2-2q-1)^2+4q(q^2-1)V.
\end{equation}
Here $c$ may be any guarded rational source function. The proof in
\cite{Trace} already permits poles of $c$ at the hypothetical bad
prime and imposes no bound on its height or denominator.

A second candidate is
\begin{equation}\label{eq:second}
                         \Lambda(c)=1+w^2V^4c^4.
\end{equation}
At a hypothetical bad pole, $\nu=v_p(w)>0$ is odd and
$v_p(V)>\nu/2$. For nonzero $c$, the valuation of $wV^2c^2$
is odd and nonzero. Therefore
\[
 v_p(\Lambda(c))=2\min\{0,v_p(wV^2c^2)\}
\]
is even; for $c=0$ it is zero. Thus $C\Lambda(c)$ has the same
pole implication as $C$. It is strictly positive on the prepared
real domain. This is an admissible sufficient test, not yet a
complete replacement for a native pole block.

For this paper, a prepared detector input means
\begin{equation}\label{eq:prepared}
 0<b<1/7,\qquad b\text{ has square reduced denominator},\qquad
 b\in\Q_2^{\times2},\quad v_2(b)\ge4,\quad v_7(b)\ge0.
\end{equation}
Then $d=7b$ meets the exact hypotheses of the intrinsic-pole
realization theorem in \cite{Intrinsic}. Our examples also satisfy
the older local square and positive-valuation requirements at
$3$ and $17$. The assertion to be tested is whether these
preparation conditions, together with a positive native product,
ensure a rational $c$ for either adaptive square equation.

\section{A universal obstruction for the first class}

\begin{lemma}\label{lem:firstlocal}
Let $p\equiv3\pmod4$ be prime and let $V\in\Q_p$ satisfy
$V\equiv1\pmod p$. Then $\Delta(q,V)$ is a nonzero square
in $\Q_p$ for every $q\in\Q_p$. Consequently, if $C$ is a
nonsquare unit, $C\Delta(q,V)$ is never a square, even with
$q$ unrestricted.
\end{lemma}
\begin{proof}
The exact identity is
\[
 \Delta(q,V)=(q^2+1)^2+4(V-1)q(q^2-1).
\]
For integral $q$, the denominator $q^2+1$ is a unit because
$-1$ is not a square modulo $p$. Thus
$\Delta/(q^2+1)^2\in1+p\Z_p$.
For $e=v_p(q)<0$, the correction quotient has valuation
$v_p(V-1)-e>0$, giving the same conclusion. The case $q=0$
gives $\Delta=1$. The subgroup $1+p\Z_p$ consists of squares
for odd $p$, proving both assertions.
\end{proof}

\begin{theorem}\label{thm:first}
There is a positive native product with
$0<w<1/4$, $200/101<V<2$ and prepared input \eqref{eq:prepared}
for which
\[
                      C\Delta(2+wc^2,V)
\]
is not a square in $\Q_{100003}$ for any $c\in\Q_{100003}$.
\end{theorem}
\begin{proof}
Put
\[
 p=100003,\qquad
 m=\frac{509796}{28707122557},\qquad
 z=\frac{509796}{32337379873},
\]
and define
\[
 u=m^2,\quad w=\frac{m^2}{(1+m^2)^2},\quad
 h=\frac{z(1+m^2)}m,\quad W=1+wh^2=1+z^2.
\]
These are rational elements of $K$, so both reciprocal slots
agree and
\[
 V=\frac2{1+z^2},\qquad
 b=\left(\frac{m(1+z^2)}{1+m^2}\right)^2.
\]
The exact checker verifies $0<m,z<1/1000$, the claimed real
bounds, and \eqref{eq:prepared}. In particular $b$ is a rational
square; its valuations at $2,3,7,17$ are all four.

Trial division through $316$ verifies that $p$ is prime, and
$p\equiv19\pmod{24}$. Direct rational reduction gives
\[
 m\equiv7121,\quad 7121^2-7121+1\equiv0,\quad z\equiv1,
 \qquad w\equiv V\equiv1,\quad C\equiv2\pmod p.
\]
Euler's criterion gives $2^{50001}\equiv-1\pmod p$.
Thus $C$ is a nonsquare unit and Lemma~\ref{lem:firstlocal}
applies to every $q$, hence to every $2+wc^2$.
\end{proof}

This obstruction is compatible with arbitrarily strong preparation
at any fixed finite set of other primes. To see this directly, choose
an integer $M$ divisible by sufficiently high powers of those primes,
with $p\nmid M$, and put
\[
                         m=\frac{M}{1+Mk},\qquad
                         z=\frac{M}{1+Ml}.
\]
Prescribe $m\equiv7121$ and $z\equiv1\pmod p$ by choosing
the classes of $k,l$ modulo $p$. Then increase positive $k,l$
in those classes. This makes $m,z$ arbitrarily small at the real
place while keeping their prescribed large positive valuations
at the fixed primes. Consequently $W\to1$, $V\to2$ there,
and $b$ is a small rational square. The radicand $C\Delta(2,V)$
can be made a nonzero local square at every fixed preparation
prime; its limit is $4\cdot49$. The fresh-prime obstruction
remains unchanged.

\section{A different obstruction for the second class}

\begin{lemma}\label{lem:secondlocal}
For $p\equiv3\pmod4$ and every $r\in\Q_p$,
\[
                         v_p(1+r^2)=2\min\{0,v_p(r)\}
\]
is even, with the usual interpretation at $r=0$. If $v_p(C)$
is odd, then $C\Lambda(c)$ is not a square for any
$c\in\Q_p$.
\end{lemma}
\begin{proof}
When $v_p(r)\ne0$ the summands have different valuations.
When $v_p(r)=0$, $-1$ is not a residue square, so the sum is
a unit. Apply this to $r=wV^2c^2$.
\end{proof}

\begin{theorem}\label{thm:second}
There is a positive native product with both real slots between
$1$ and $101/100$, with normalized weight and prepared input
\eqref{eq:prepared}, such that $C\Lambda(c)$ is not a square
in $\Q_{100003}$ for any $c\in\Q_{100003}$.
\end{theorem}
\begin{proof}
Keep $p=100003$ and set
\[
 r=44974,\qquad M=714^2,\qquad
 h=Mp^3(r+\sqrt2),
\]
\[
 u=52581464203154325228257242401217852539664811276478096.
\]
\[
 w=\frac{u}{(1+u)^2},\qquad W=1+wh^2.
\]
Write $W=A+B\sqrt2$, $N=A^2-2B^2$, and use
$V=2A/N$, $b=4w/V^2$. The integer $u$ satisfies
\[
 \begin{array}{c|ccccc}
 \text{modulus}&p^3&128&27&4913&7\\\hline
 u& p^2-1&16&9&289&1.
 \end{array}
\]
It was chosen by the Chinese remainder theorem above
$100H^2+100$, where $H=Mp^3(r+2)$ bounds the absolute
value of both real slots of $h$. Since $w<1/u$, both slots
of $W$ lie in $(1,101/100)$. Hence $200/101<V<2$,
and the exact bound also gives $0<b<1/7$.

All odd valuations of $w$ are positive because $u$ is a positive
integer and its normalization denominator is a square. Moreover
\begin{equation}\label{eq:bsquareclass}
                         b=u\left(\frac{N}{(1+u)A}\right)^2.
\end{equation}
At every odd prime dividing $u$, the integrality of $h$ gives
$W\equiv1$, $A,N\equiv1$, and $V\equiv2$. Thus $b$ has
positive valuation there. At all other odd primes its valuation
is even. The displayed congruences give $v_2(b)=4$ and the
$2$-adic square class of a square. They also give the square
classes and valuations two at $3$ and $17$. At $7$, $h$ is
divisible by $7$ and $u\equiv1$, so $A,N,V$ are units and
$v_7(b)=0$. This proves \eqref{eq:prepared}, including square
reduced denominator. The checker also verifies that denominator
directly as an integer square.

At the fresh prime, $v_p(r^2+2)=1$ and $v_p(1+u)=2$.
Both slots of $h$ have valuation three in the unramified
quadratic extension. With $t=wh^2$ one has $v_p(t)=2$ and
\[
 v_p(\Tr(t))=v_p\bigl(2wM^2p^6(r^2+2)\bigr)=3.
\]
The terms from $t^2,t^3,\ldots$ in the convergent expansion of
$(1+t)^{-1}$ have trace valuation at least four. Consequently
\[
 v_p(w)=-4,\quad v_p(V)=0,\quad v_p(2-V)=3,
 \quad v_p(C)=-1,\quad v_p(b)=-4.
\]
The negative valuation of $b$ is even and is compatible with
pole preparation. The odd valuation of $C$ invokes
Lemma~\ref{lem:secondlocal}, excluding every scalar $c$.
\end{proof}

Here too any prescribed finite preparation away from $p$ can
be strengthened. Increase the divisibility of $M$ at the fixed
primes. At each of them choose $u$ with arbitrarily large even
positive valuation and square unit part, while retaining
$u\equiv p^2-1\pmod{p^3}$. Increase the positive CRT solution
so that $u/H^2$ is arbitrarily large. Then $W$ approaches one
and $V$ approaches two at the fixed places and at the real place,
while $b$ remains prepared. The proof of the fresh valuations
uses only the unchanged congruences at $p$.

\section{The quartic curve behind the second criterion}

Let $w>0$ and $1<V<2$ be fixed rational values. Put
\[
 B_*=\frac C{V^2}=1+\frac b4(2-V),\qquad z=\frac{V^2c}{2}.
\]
Dividing the square root by $V$ transforms the second criterion
exactly into
\begin{equation}\label{eq:quartic}
                         y^2=B_*(1+b^2z^4).
\end{equation}
Its nonsingular projective normalization has genus one. For
$z\ne0$, the explicit map
\[
 X=B_*b^2z^2,\qquad Y=B_*b^2zy
\]
gives
\[
                         Y^2=X^3+(B_*b)^2X.
\]
Conversely a point on this latter curve with
$X\in B_*\Q^{\times2}$ gives a point of \eqref{eq:quartic}.
The point $z=0$ exists exactly when $B_*$ is a rational square.
Writing $n=B_*b$, the usual degree-two map
\[
 x=X+\frac{n^2}{X},\qquad y'=Y\left(1-\frac{n^2}{X^2}\right)
\]
lands on the prescribed congruent-number curve
\[
                         (y')^2=x^3-(2n)^2x.
\]
The checker verifies both maps by exact polynomial identities.
They preserve an additional square-class requirement; mere positive
rank of the last curve would not automatically recover the
prescribed quartic point.

The intrinsic-pole theorem permits choosing an auxiliary twist
parameter after the prepared input has been fixed. It does not
prescribe a point on this fixed curve with parameter $2B_*b$.
Theorem~\ref{thm:second} directly shows that pole preparedness
cannot supply such a conclusion in general.

For comparison, the first criterion is already a higher-genus
problem in its unrestricted scalar $c$. Its quartic in $q$ has
discriminant
\[
                 16384V^2(V-2)^2(V-1)^2.
\]
For $1<V<2$ it has four distinct roots, none equal to $2$,
since $\Delta(2,V)=1+24V>0$. Substituting $q=2+wc^2$
therefore gives a squarefree degree-eight polynomial, whose
double cover has genus three. This observation is independent
of the two exact local obstructions.

\section{What the counterexamples do and do not establish}

Both examples are literal positive native products of the form
$1+wh^2$ over $K$. They satisfy the normalized-weight parity
condition, the intrinsic pole preparation, and strict real bounds;
arbitrary fixed-place preparation away from the obstruction prime
is compatible with the obstructions.
Under the retained arithmetic premise, the intrinsic-pole theorem
still supplies its two-parameter witnesses at $d=7b$.

The examples therefore disprove completeness of either proposed
one-free-scalar replacement on the class of all such prepared
native inputs. The first obstruction permits every $q\in\Q_p$,
which is stronger than merely varying $c$ in its specified formula.
The second obstruction defeats every rational value of $c$ in
\eqref{eq:second}.

Neither $h$ nor $u$ in these examples has been lifted to the actual
elliptic test tuple and physical contact map of the general or
conditional source. These results are not counterexamples to
existence of some other actual prepared witness for each integer
input. No impossibility of a complete twenty-one- or nine-variable
definition, and no lower bound on defining-polynomial variables,
is asserted.

The standalone exact verifier and full rational values are in
\begin{quote}\small
\nolinkurl{negative_results/computation/pole_parameters/adaptive_prepared_pairs_v76}.
\end{quote}
The universal quantifiers over $c$ follow from the valuation
lemmas, not from a finite search. The checker verifies primality,
residue and valuation data, reduced denominators, all stated
preparation inequalities, and the polynomial identities.

\begin{thebibliography}{9}\raggedright
\bibitem{Trace}\emph{A rational conic criterion using the actual
detector trace},
\nolinkurl{papers/pole_parameters/trace_conic_criterion_20261117_v57.tex},
especially the adaptive-class proposition in its final subsection.
\bibitem{Intrinsic}\emph{A two-coordinate pole test of coefficient
height four},
\nolinkurl{papers/pole_parameters/intrinsic_pole_20261117_v56.tex},
Theorem 2.1 and its prepared detector-input application.
\end{thebibliography}
\end{document}
